% ============================================================================
%  Hessian-aware weight reconstruction is conditional.
%
%  Standalone-compilable (pdflatex hessian_conditional.tex), or strip the
%  preamble and \input the section body into the paper.
%
%  CLAIM (conditional / honest): the Hessian-damped, error-feedback weight
%  calibration (Lloyd/GPTQ-style, alpha* = (b^T H w)/(b^T H b), damped) is
%  NECESSARY only when the layer's activation Hessian H = X^T X / T retains an
%  extreme, quantization-axis-aligned spike.  The randomized Hadamard rotation
%  that removes activation outliers ALSO de-aligns the heavy eigendirections of
%  H from the quantization axes, shrinking the RTN-vs-optimal gap by ~sqrt(d);
%  for generic layers this makes plain round-to-nearest (RTN) near-optimal
%  ("Hessian barely helps single-GF4"), while pathological layers (opt-13b
%  W4A16) still collapse under RTN (304) and require the reconstruction (10.5).
%
%  Grounding (CUDA_FP4_Test/reference.py, verified):
%    H_block = X_block^T X_block / T          (bit_split.py)
%    alpha*  = (b_eff . (H w)) / (b_eff . (H b))     Lloyd step, per block
%    damping lam_add = max(lam_max/kappa - lam_min, 0),  kappa = 100
%    opt-13b W4A16 collapse 304 -> 10.5 with the damped reconstruction.
% ============================================================================
\documentclass[11pt]{article}
\usepackage{amsmath,amssymb,amsthm,bm}
\usepackage[margin=1in]{geometry}
\usepackage{booktabs}

\newtheorem{lemma}{Lemma}
\newtheorem{theorem}{Theorem}
\newtheorem{corollary}{Corollary}
\theoremstyle{remark}
\newtheorem{remark}{Remark}
\newtheorem{assumption}{Assumption}

\newcommand{\R}{\mathbb{R}}
\newcommand{\E}{\mathbb{E}}
\newcommand{\tr}{\mathrm{tr}}
\newcommand{\norm}[1]{\lVert #1 \rVert}
\newcommand{\Hm}{\bm H}
\newcommand{\wh}{\hat{\bm w}}

\begin{document}

\section{Hessian-aware reconstruction is conditional}
\label{sec:hessian}

\subsection{Setup: the layer objective and what RTN ignores}

For a linear layer $\bm y=\bm x\bm W^{\!\top}$, the calibration-optimal
quantized weight minimizes the expected \emph{output} error, not the weight
error:
\begin{equation}
\label{eq:obj}
\min_{\hat{\bm W}\in\mathcal Q}\;
\E_{\bm x}\big\|(\bm W-\hat{\bm W})\bm x\big\|^2
= \min_{\hat{\bm W}\in\mathcal Q}\;
\tr\!\big[(\bm W-\hat{\bm W})\,\Hm\,(\bm W-\hat{\bm W})^{\!\top}\big],
\qquad
\Hm=\E[\bm x\bm x^{\!\top}]\approx\tfrac1T \bm X^{\!\top}\bm X,
\end{equation}
where $\Hm\in\R^{d\times d}$ is the (per-layer) activation second moment --- the
Gauss--Newton Hessian of the reconstruction. The objective decouples over rows;
for a single row $\bm w\in\R^d$ with quantization residual
$\bm\delta=\bm w-\wh$,
\begin{equation}
\label{eq:E}
E(\wh)=\bm\delta^{\!\top}\Hm\,\bm\delta
      =\underbrace{\textstyle\sum_i \Hm_{ii}\delta_i^2}_{\text{diagonal / unavoidable}}
      +\underbrace{\textstyle\sum_{i\neq j}\Hm_{ij}\delta_i\delta_j}_{\text{cross terms}} .
\end{equation}
\textbf{Round-to-nearest (RTN)} quantizes each coordinate independently, i.e.\
minimizes $\norm{\bm\delta}^2$ ($\Hm{=}\bm I$); it ignores the second term.
\textbf{Hessian-aware reconstruction} (the damped Lloyd/GPTQ step
$\alpha^\star=(\bm b^{\!\top}\Hm\bm w)/(\bm b^{\!\top}\Hm\bm b)$ with error
feedback) instead \emph{steers} $\bm\delta$ into the low-curvature subspace to
cancel the cross terms. The question is when that steering is worth its cost.

\subsection{Lemma 1: the RTN gap is set by curvature aligned with the axes}

\begin{lemma}[RTN suboptimality]
\label{lem:gap}
Let $\bm\delta$ be the RTN residual (per coordinate, $\E\delta_i=0$,
$\E\delta_i^2=\sigma^2=\Delta^2/12$, independent across $i$). Then
\begin{equation}
\label{eq:gapmean}
\E\,E_{\mathrm{RTN}} = \sigma^2\tr(\Hm)
= \sigma^2\!\sum_i\Hm_{ii},
\qquad
\E\,E_{\mathrm{RTN}}-E_{\mathrm{opt}}
\;\le\; \sigma^2\big(\tr\Hm - d\,\lambda_{\min}(\Hm)\big),
\end{equation}
and the achievable reduction $E_{\mathrm{RTN}}-E_{\mathrm{opt}}$ is nonzero only
to the extent the residual can be rotated toward small-eigenvalue directions of
$\Hm$; it scales with the \emph{anisotropy} $\lambda_{\max}/\bar\lambda$
($\bar\lambda=\tr\Hm/d$) \emph{as seen in the quantization (coordinate) basis}.
\end{lemma}

\begin{proof}
Independence and zero mean kill the cross terms in expectation, giving
$\E E_{\mathrm{RTN}}=\sigma^2\sum_i\Hm_{ii}=\sigma^2\tr\Hm$. Any admissible
quantizer pays at least $\sigma^2 d\,\lambda_{\min}$ on the residual energy it
cannot remove, so the optimal value is bounded below by that; subtracting gives
\eqref{eq:gapmean}. The reduction is realized only by choosing $\bm\delta$
correlated with the eigenvectors of $\Hm$ (error feedback), whose leverage is
the eigenvalue spread.
\end{proof}

The gap is large exactly when a \emph{few} coordinates carry large curvature
$\Hm_{ii}$ (a spike aligned with an axis): RTN spends $\sigma^2\Hm_{ii}$ there,
while reconstruction keeps those coordinates precise and dumps the error into
flat directions. This is the regime where GPTQ/OBQ-style calibration earns its
keep.

\subsection{Lemma 2: rotation de-aligns the spike from the quantization axes}

Let $\bm R=\tfrac1{\sqrt d}\bm H\bm D$ be the randomized Hadamard transform used
for the activations; the layer is quantized in the rotated basis, with Hessian
$\tilde{\Hm}=\bm R\,\Hm\,\bm R^{\!\top}$. A similarity transform preserves the
spectrum, so rotation does \emph{not} change $\{\lambda_i\}$ --- but it changes
how that spectrum sits relative to the (axis-aligned) quantization lattice.

\begin{lemma}[Incoherence of the rotated Hessian]
\label{lem:rotH}
Fix $\Hm\succeq 0$. For $\bm R$ as above, with probability $\ge 1-\delta$,
\begin{equation}
\label{eq:diag}
\max_i\big|\tilde{\Hm}_{ii}-\bar\lambda\big|
   \;\le\; \frac{\lambda_{\max}}{\sqrt d}\,\sqrt{2\log\tfrac{2d}{\delta}},
\qquad
\max_{i\neq j}\big|\tilde{\Hm}_{ij}\big|
   \;\lesssim\; \frac{\lambda_{\max}}{\sqrt d}\,\sqrt{2\log\tfrac{2d}{\delta}},
\end{equation}
i.e.\ every diagonal entry concentrates at $\bar\lambda=\tr\Hm/d$ and every
off-diagonal is $O(\lambda_{\max}/\sqrt d)$. Consequently any single
eigendirection $\bm v$ of $\Hm$ has coordinates
$|(\bm R\bm v)_i|=O(1/\sqrt d)$ in the rotated basis (delocalization).
\end{lemma}

\begin{proof}
$\tilde{\Hm}_{ij}=\bm r_i^{\!\top}\Hm\,\bm r_j$ with $\bm r_i=\bm R^{\!\top}\bm e_i$
a randomized-sign unit-norm row; write $\Hm=\sum_k\lambda_k\bm u_k\bm u_k^{\!\top}$.
Each $\langle\bm r_i,\bm u_k\rangle$ is a Rademacher sum with sub-Gaussian proxy
$1/d$, so $\bm r_i^{\!\top}\Hm\bm r_j$ concentrates at
$\bar\lambda\,\mathbf 1\{i{=}j\}$ with fluctuation
$O(\lambda_{\max}/\sqrt d)$ by Hanson--Wright / a union bound over the $d^2$
entries; this is \eqref{eq:diag}. Delocalization is the same sub-Gaussian bound
applied to $(\bm R\bm v)_i=\bm r_i^{\!\top}\bm v$.
\end{proof}

The point is \emph{not} that rotation whitens $\Hm$ (it cannot; the eigenvalues
are fixed) but that it makes $\tilde{\Hm}$ \emph{diagonally dominant with a
uniform diagonal}: the heavy curvature no longer aligns with any coordinate. By
Lemma~\ref{lem:gap} the exploitable gap is governed by axis-aligned curvature,
which \eqref{eq:diag} caps at $\bar\lambda+O(\lambda_{\max}/\sqrt d)$.

\subsection{Theorem: when reconstruction is (un)necessary}

\begin{theorem}[Conditional necessity of Hessian reconstruction]
\label{thm:cond}
In the rotated basis, the RTN suboptimality obeys
\begin{equation}
\label{eq:cond}
0\;\le\;\E\,E_{\mathrm{RTN}}-E_{\mathrm{opt}}
   \;\le\; \sigma^2\,\frac{\lambda_{\max}}{\sqrt d}\,
           \sqrt{2\log\tfrac{2d}{\delta}}\;\cdot\;\rho_{\mathrm{off}},
\end{equation}
where $\rho_{\mathrm{off}}\le d$ counts the residual off-diagonal coupling.
Hence:
\begin{itemize}
\item[(i)] If the layer is well conditioned or only mildly spiked after
rotation, $\lambda_{\max}=O(\bar\lambda\sqrt d)$, the bound is
$O(\sigma^2\bar\lambda)$ --- the same order as the unavoidable diagonal term
--- so RTN is near-optimal and \textbf{Hessian reconstruction is unnecessary}.
\item[(ii)] If a residual extreme spike survives rotation,
$\lambda_{\max}\gtrsim\bar\lambda\,d$ (an $\Omega(\sqrt d)$ outlier eigenvalue
that Hadamard mixing cannot equalize), the bound is
$\Omega(\sigma^2\bar\lambda\sqrt d)$ and RTN can diverge; \textbf{the damped
reconstruction is necessary} and recovers the $O(\sigma^2\bar\lambda)$ optimum.
\end{itemize}
\end{theorem}

\begin{proof}
Combine Lemma~\ref{lem:gap} (gap $\le\sigma^2(\tr\tilde\Hm-d\lambda_{\min})$ and
proportional to axis-aligned curvature) with Lemma~\ref{lem:rotH} (axis-aligned
curvature $\le\bar\lambda+O(\lambda_{\max}/\sqrt d)$), yielding \eqref{eq:cond}.
Cases (i)/(ii) substitute the two scalings of $\lambda_{\max}$; in case (ii) the
$\alpha^\star$-Lloyd step attains $E_{\mathrm{opt}}$ up to the damping bias of
Remark~\ref{rem:damp}.
\end{proof}

\begin{corollary}[The rotation is doing the calibration's job --- until it can't]
The same randomized Hadamard rotation that removes activation outliers
(\S\ref{sec:retention}, Lemma on incoherence) simultaneously de-aligns the
weight Hessian's heavy directions, so on \emph{most} layers RTN already sits at
the Hessian optimum. Reconstruction is therefore a targeted fix for the few
layers whose curvature spike is too large for $\sqrt d$ mixing to absorb, not a
uniform requirement.
\end{corollary}

\subsection{The damping term, and why it is the right conditional knob}

\begin{remark}[Damping]
\label{rem:damp}
The reconstruction solves $\alpha^\star=(\bm b^{\!\top}\Hm\bm w)/
(\bm b^{\!\top}\Hm\bm b)$ on a Hessian regularized by
$\Hm\leftarrow\Hm+\lambda_{\mathrm{add}}\bm I$,
$\lambda_{\mathrm{add}}=\max(\lambda_{\max}/\kappa-\lambda_{\min},0)$,
$\kappa=100$. This caps the condition number of the working Hessian at $\kappa$:
it is active precisely in case (ii) (a genuine spike, $\lambda_{\min}\ll
\lambda_{\max}/\kappa$) and inert in case (i) ($\lambda_{\min}\ge\lambda_{\max}/
\kappa\Rightarrow\lambda_{\mathrm{add}}=0$, reconstruction $\to$ RTN). Damping is
thus the mechanism that makes the procedure \emph{degrade gracefully to RTN}
exactly where Theorem~\ref{thm:cond}(i) says RTN suffices.
\end{remark}

\subsection{Empirical instantiation}

\begin{center}\small
\begin{tabular}{@{}lccl@{}}
\toprule
Layer / model regime & RTN (H-agnostic) & Damped reconstruction & Verdict \\
\midrule
opt-13b, W4A16, padded MLP (spiked $\Hm$)
   & $\approx\!304$ ppl (collapse) & $\approx\!10.5$ ppl & reconstruction \emph{necessary} \\
generic layers, single-GF4 W4A4
   & near-optimal & $\approx$ same (barely helps) & reconstruction \emph{optional} \\
\bottomrule
\end{tabular}
\end{center}

The two rows are the two branches of Theorem~\ref{thm:cond}. The opt-13b
padded-MLP layer carries an $\Omega(\sqrt d)$ curvature spike that survives
rotation (case ii): RTN collapses to $304$, the damped $\alpha^\star$-Lloyd step
restores $10.5$. On the remaining layers the rotated Hessian is diagonally
dominant (case i): RTN is already at the optimum, so adding the reconstruction
moves perplexity within noise --- consistent with the observation that a single
Hessian pass ``barely helps'' the generic GF4 build.

\subsection{What is and is not claimed}

\begin{remark}[Scope]
The claim is \emph{conditional}: Hessian-aware reconstruction is not a blanket
requirement for FP4 weight quantization. Given the activation rotation, plain
RTN is provably near-optimal for any layer whose post-rotation curvature spike
is $O(\bar\lambda\sqrt d)$ (the common case); reconstruction is required only for
residual extreme spikes ($\lambda_{\max}=\Omega(\bar\lambda d)$), which in our
suite means specifically the pathological padded/outlier layers (opt-13b). The
damping $\kappa=100$ makes the method collapse to RTN automatically where it is
not needed.
\end{remark}

\begin{remark}[Honest caveats]
(a) Lemma~\ref{lem:rotH} is high-probability over the rotation, not worst-case;
an adversarial $\Hm$ co-designed with $\bm R$ is excluded. (b) The residual model
$\E\delta_i\delta_j=\sigma^2\mathbf 1\{i{=}j\}$ (Lemma~\ref{lem:gap}) is the
standard high-resolution surrogate; real integer residuals are weakly correlated
and only rescale constants. (c) $E_{\mathrm{opt}}$ is the coordinate-descent
(GPTQ) optimum, not the NP-hard lattice optimum; the $\alpha^\star$ step attains
it up to the damping bias. (d) The empirical row values are single-layer /
whole-model perplexities under the paper's protocol and inherit its eval noise.
\end{remark}

\end{document}
